\documentclass[11pt,a4paper]{article}

\usepackage[margin=2.5cm]{geometry}
\usepackage[T1]{fontenc}
\usepackage[utf8]{inputenc}
\usepackage{lmodern}
\usepackage{microtype}
\usepackage{booktabs}
\usepackage{tabularx}
\usepackage{longtable}
\usepackage{array}
\usepackage{enumitem}
\usepackage{xcolor}
\usepackage{hyperref}
\usepackage{caption}
\usepackage{float}

\definecolor{smublue}{RGB}{21,55,89}
\hypersetup{
  colorlinks=true,
  linkcolor=smublue,
  urlcolor=smublue,
  citecolor=smublue,
  pdftitle={From Forum Posts to Telco Actions},
  pdfauthor={Zixin}
}
\setlength{\parindent}{0pt}
\setlength{\parskip}{0.55em}
\setlist{nosep}
\renewcommand{\arraystretch}{1.18}

\newcommand{\studentname}{Zixin [Surname]}
\newcommand{\studentid}{[Student ID]}
\newcommand{\programme}{[Programme / Course]}
\newcommand{\supervisorname}{[Professor's Name]}

\title{\textbf{From Forum Posts to Telco Actions}\\
\large An End-to-End Pipeline for Sentiment and Issue Analysis\\
of Singapore Telecommunications Discussions}
\author{\studentname\\\studentid\\\programme}
\date{September 2026}

\begin{document}

\begin{center}
\fcolorbox{black}{gray!12}{\parbox{0.94\textwidth}{\centering
\textbf{SMU Classification: Restricted}}}
\end{center}

\maketitle

\begin{center}
\textbf{Supervisor:} \supervisorname
\end{center}

\vfill

\begin{abstract}
Singapore telecommunications operators receive continuous, unsolicited customer feedback through online forums, but the volume, informality and platform-specific language make manual monitoring difficult. This project develops an end-to-end pipeline that collects forum discussions, normalises heterogeneous records, identifies telecommunications brands, assigns sentiment and issue labels, and presents findings through a static HTML dashboard. The analysis corpus combines Reddit and HardwareZone (HWZ) posts up to 26~September~2026: after removing repeated copies and off-topic threads, 9,162 unique posts in 349 telco threads remain. GPT-4o-mini labels 44\% of them not happy and 19\% happy (42\% and 18\% with threads weighted equally), whereas VADER on the same posts reports 28\% and 46\%. Network quality is the largest issue theme (37--41\% of posts), followed by mobile plans (13--14\%) and pricing (7--8\%). Reddit and HWZ differ mainly in topic rather than tone, and no brand stands out beyond small, weighting-sensitive differences. One event shows why that care is needed: 45 outage threads, led by the March 2026 Singtel outage, hold 12\% of posts, are 71\% not happy against 40\% elsewhere, and sit almost entirely on Reddit, which creates most of the apparent Reddit--HWZ gap.

The system combines deterministic brand matching, VADER sentiment, a keyword issue taxonomy, an optional Hugging Face sentiment model and OpenAI chat-completions classification. Independent human labels on a blinded 300-row representative sample (295 telco-relevant rows) show that \textbf{gpt-4o-mini-2024-07-18} improves sentiment over VADER (macro F1 0.631 versus 0.360; accuracy 67\% versus 37\%). \textbf{gpt-4o-2024-08-06} is similar on sentiment (macro F1 0.623) and still loses to keyword issue rules when it labels every row (macro F1 0.600 versus 0.639). A hybrid that keeps rule labels and applies GPT-4o only to unresolved \texttt{general\_discussion} rows raises issue macro F1 to \textbf{0.680}, and was applied to all 12,497 unresolved corpus rows (about 41 minutes, estimated \textbf{\$4.80}). On a held-out 100-post challenge set of harder posts, the sentiment result holds (GPT-4o-mini 0.773 versus VADER 0.491), but the hybrid's gain over rules is not significant (0.509 versus 0.499) and GPT-4o labelling every post scores best (0.596), so no single issue method wins on both sets. GPT-4o-mini sentiment on all 9,162 posts cost about \$0.28, and a repeated run reproduced 292 of 295 gold-set labels. On a third 100-post gold set drawn from posts published between April and September 2026, GPT-4o-mini scores macro F1 0.637 against VADER's 0.428, so the sentiment result carries over, although the model under-calls complaints there, so the corpus not-happy share is likely an underestimate. An earlier 150-record VADER ``100\% accuracy'' file is discarded because its labels match VADER exactly and were model-prefilled. The operational recommendation is therefore LLM sentiment (mini is sufficient), keyword rules plus a targeted GPT-4o pass for issues, and human review of alerts. The business analysis indicates that operators should prioritise incident communication during outages, comparable plan terms such as fair-use limits and add-on charges, and equipment and app friction, with billing and customer-service aftercare planned for incidents; roaming and porting are small in this corpus. The principal contribution is a reproducible analytics workflow that converts noisy forum data into decision-oriented signals while making its evidential limits explicit.
\end{abstract}

\vfill
\begin{center}
\small\textit{Keywords:} telecommunications, social listening, sentiment analysis, issue extraction, VADER, large language model, dashboard
\end{center}
\newpage

\tableofcontents
\newpage

\section{Introduction}

\subsection{Business context}

Mobile customers discuss price plans, coverage, roaming, porting and service failures in public online communities. These discussions can reveal emerging pain points earlier and in more detail than periodic surveys. They are also difficult to analyse reliably: posts contain Singlish, abbreviations, sarcasm, quoted text, multiple brands and off-topic uses of ambiguous brand names such as ``Eight'' or ``M1''.

This project asks how a telecommunications operator can convert such unstructured discussion into a repeatable source of customer insight. The objective is not to infer population-level satisfaction or replace operational metrics. Instead, the pipeline is designed as an early-warning and exploratory decision-support system that helps analysts locate themes, compare discussion contexts and inspect representative records.

\subsection{Objectives}

The project has five objectives:

\begin{enumerate}
  \item collect relevant discussions from multiple online platforms;
  \item clean and standardise heterogeneous post and comment records;
  \item detect mentioned operators and classify sentiment and service issues;
  \item compare lightweight baselines with more context-aware model options; and
  \item translate the results into actionable recommendations through a dashboard.
\end{enumerate}

\subsection{Research questions}

\begin{description}[leftmargin=0cm,style=nextline]
  \item[RQ1.] What sentiment and issue themes appear most frequently in the collected Singapore telco discussions?
  \item[RQ2.] How do discussion patterns differ between Reddit and HWZ, and across telco brands?
  \item[RQ3.] What additional value can a context-aware model provide over VADER sentiment and keyword-based issue rules?
  \item[RQ4.] How can the extracted signals support concrete decisions by telco operators?
\end{description}

\subsection{Scope and claimed contribution}

The repository implements collection modules for HWZ and Reddit. It includes preprocessing, brand detection, sentiment labelling, issue extraction, evaluation scripts and a static HTML dashboard. The report only claims findings supported by persisted repository outputs. GPT-4o-mini and GPT-4o are evaluated on human labels for both the representative and challenge sets, and GPT-4o is used on unresolved corpus issues. The findings and dashboard use GPT-4o-mini sentiment and hybrid issue labels on the 9,162-post corpus.

\section{Data and Pipeline}

\subsection{Data sources}

The corpus covers Reddit posts from Singapore-focused communities, matched on telco and mobile-plan keywords, and HWZ posts from relevant XenForo threads, up to 26~September~2026. Reddit posts up to 9~April~2026 came from the public JSON interface; because Reddit then blocked unauthenticated access and gated new API applications, later Reddit posts and comments (including replies) came from the Arctic Shift research archive, filtered with the same strict brand rule. A page-counter bug meant early HWZ scrapes read only the first page of each thread; after the fix, the newest pages of each telco thread were collected as well. The scrapers preserve thread titles, post or comment text, timestamps, source identifiers and available engagement metadata.

\begin{table}[H]
\centering
\caption{Persisted analytical datasets}
\label{tab:data}
\begin{tabularx}{\textwidth}{@{}p{0.34\textwidth}rX@{}}
\toprule
\textbf{Dataset} & \textbf{Rows} & \textbf{Role} \\
\midrule
\texttt{issues\_posts\_labeled\_\ldots\_gpt4o\_hybrid.csv} & 24,690 & All scraped rows after preprocessing, with brand, VADER and issue labels (rules plus GPT-4o on 12,497 unresolved rows) \\
\texttt{corpus\_posts.csv} & 9,162 & Clean corpus: unique posts in 349 telco threads \\
\texttt{corpus\_posts\_mini\_sentiment.csv} & 9,162 & Clean corpus with GPT-4o-mini sentiment; findings and dashboard input \\
\texttt{gold\_labels.csv} & 500 & Annotator~A's three blind gold sets in one file, scored separately by \texttt{evaluation\_set}: representative (300; 295 telco-relevant, headline metrics), challenge (100 harder posts) and recent (100 posts from April--September 2026; 97 relevant) \\
\texttt{llm\_comparison\_representative/} & 295 & GPT-4o-mini vs VADER/rules predictions, metrics, tokens and cost \\
\texttt{llm\_comparison\_representative\_gpt4o/} & 295 & GPT-4o vs VADER/rules on the same gold labels \\
\texttt{llm\_comparison\_representative\_gpt4o\_hybrid/} & 295 & Rules, then GPT-4o only where rules output \texttt{general\_discussion} \\
\texttt{llm\_comparison\_challenge\_gpt4o/} & 100 & GPT-4o vs VADER/rules on the challenge set; hybrid scored from it \\
\texttt{robustness\_summary.json} & --- & Bootstrap intervals, run-to-run stability, annotator agreement, corpus checks \\
\texttt{labeled\_sample.csv} & 150 & Earlier Reddit sample; labels not independent; superseded \\
\texttt{hwz\_telco\_posts\_labeled.csv} & 100 & Earlier HWZ sample for the 158-post dashboard; superseded \\
\bottomrule
\end{tabularx}
\end{table}

The scraped rows are not independent posts. The 24,690 labelled rows hold 11,643 unique posts, because successive scrape rounds saved the same comments again; in one ten-round scrape, rounds 2--9 added no new posts. Ninety-six of 445 threads were off-topic and matched a keyword by accident, such as a 248-post thread on what feels cheaper in Singapore and a 98-post insurance lawsuit thread. \texttt{processing/build\_clean\_corpus.py} keeps one row per post ID (Reddit comment or post ID; HWZ thread URL and timestamp), so identical text from different users stays separate, and drops a thread unless at least 20\% of its posts name a telco or its title is telco-related. The result is 9,162 posts in 349 threads (5,711 Reddit, 3,451 HWZ). An earlier 158-post dashboard sample (67 Reddit, 91 HWZ) is superseded.

\subsection{Preprocessing}

The preprocessing module maps platform-specific fields to a common schema, removes empty, placeholder and very short records, and can optionally retain English-language records using language detection. The common schema includes \texttt{source}, \texttt{thread\_title}, \texttt{post\_text}, timestamp fields and identifiers. Normalisation allows downstream modules to process different platforms without source-specific model code.

The current workflow retains post text that may include informal spelling, emoji and HTML entities because these features can carry sentiment. This is useful for analysis but also creates modelling challenges. The dashboard therefore presents aggregate indicators and traceable examples rather than treating model output as verified customer facts.

\subsection{End-to-end architecture}

The implemented sequence is:

\begin{center}
\begin{minipage}{0.92\textwidth}
\small
\textbf{Collection} (Reddit and HWZ scrapers)
$\rightarrow$ \textbf{raw CSV}
$\rightarrow$ \textbf{preprocessing}
$\rightarrow$ \textbf{brand and sentiment labelling}
$\rightarrow$ \textbf{issue extraction}
$\rightarrow$ \textbf{evaluation outputs}
$\rightarrow$ \textbf{static HTML dashboard}.
\end{minipage}
\end{center}

The static dashboard generator is deliberately lightweight and produces a standalone HTML report that opens in any browser without a server.

\section{Analytical Methods}

\subsection{Telco detection}

Brand mentions are identified by regular expressions over the concatenated thread title and post text. The tracked labels include Singtel, StarHub, M1, SIMBA, Giga, Eight, Zym and redONE, with additional MVNO keywords used by the dashboard. Multiple brands can be assigned to one record. This design is transparent and inexpensive, but ambiguous names and thread-level context can create false positives. For example, a post containing the ordinary word ``eight'' may be incorrectly associated with the Eight brand unless the surrounding context is considered.

\subsection{Sentiment baseline: VADER}

VADER is a lexicon- and rule-based sentiment method designed for social-media text \cite{hutto2014vader}. For each post, the pipeline obtains a compound score and maps it to:

\begin{equation}
\text{label}(s)=
\begin{cases}
\text{happy}, & s \geq 0.05,\\
\text{not\_happy}, & s \leq -0.05,\\
\text{neutral}, & \text{otherwise}.
\end{cases}
\end{equation}

VADER is a suitable operational baseline because it runs locally, is deterministic and requires no paid API. Its limitation is contextual understanding. It can misread sarcasm, negation over long passages, mixed opinions, quoted complaints and the difference between sentiment towards a telco and sentiment towards another topic in the same thread.

\subsection{Context-aware sentiment option}

The benchmark module supports \texttt{cardiffnlp/twitter-roberta-base-sentiment-latest}, a RoBERTa-family model adapted to social-media sentiment \cite{loureiro2022tweeteval}. Its positive, negative and neutral outputs are mapped to the project's three business labels. The persisted benchmark states that the required \texttt{transformers} package was unavailable, so no transformer metrics are reported.

The main labelling module also implements LLM sentiment through an OpenAI-compatible chat-completions interface. Its configured default is \textbf{GPT-4o-mini}. Formal comparison uses the snapshots returned by the API, \textbf{gpt-4o-mini-2024-07-18} and \textbf{gpt-4o-2024-08-06}, with temperature 0 and JSON object responses in batches of ten posts. An earlier 100-record HWZ file contains 100 neutral LLM labels and no model-version metadata; that file is treated as a failed or fallback run and is not used as evidence. The corpus was labelled with the same prompt, batching and settings through \texttt{models/corpus\_llm\_sentiment.py}; the checkpoint is keyed by post ID, so posts already labelled in an earlier run were not sent again. Labelling all 9,162 posts cost about \$0.28.

\subsection{Issue taxonomy and rationale}

The issue taxonomy contains ten business-facing categories. These labels were selected because they separate product, network, service, billing, digital-channel and device decisions while retaining a general-discussion fallback:

\begin{table}[H]
\centering
\caption{Issue taxonomy and business rationale}
\label{tab:taxonomy}
\begin{tabularx}{\textwidth}{@{}p{0.19\textwidth}p{0.35\textwidth}X@{}}
\toprule
\textbf{Label} & \textbf{Typical evidence} & \textbf{Action owner / decision} \\
\midrule
Network Quality & coverage, signal, 4G/5G, speed, latency, outage & Network teams: local investigation, capacity and coverage communication \\
Pricing \& Promotions & price, promotion, discount, deal, cashback & Commercial teams: price communication and promotion fulfilment \\
Mobile Plans & SIM-only allowances, contract terms, minutes, SMS & Product teams: plan design and feature comparison \\
Roaming & overseas data, IDD, APAC, international use & Roaming teams: bundles, activation and destination guidance \\
Porting & port-in/out, number transfer, switching & Operations: mobile-number-portability workflow and exception handling \\
Customer Service & support, hotline, agent, complaint, escalation & Service operations: contact journeys and case resolution \\
Billing \& Payment & bills, charges, refunds, invoices, GIRO & Billing teams: charging controls and payment resolution \\
App \& Account & app, login, password, OTP, verification, portal & Digital teams: self-service access and account management \\
Devices & phones, eSIM/SIM, routers, ONR/ONT, equipment & Device and technical teams: compatibility and equipment support \\
General Discussion & telco news, opinions, no specific issue & Analysts: monitor broad conversation and review taxonomy coverage \\
\bottomrule
\end{tabularx}
\end{table}

The categories remain compact enough to route to owners while separating previously conflated topics such as plan design from promotions, billing from customer service, and app access from device problems. Multi-label output permits a post to mention, for example, both pricing and roaming. The rule-based classifier assigns labels when predefined regular expressions match and otherwise assigns \texttt{general\_discussion}. GPT-4o was then applied only to those unresolved rows. On the representative gold set that hybrid beats both keyword-only and all-row LLM tagging, so it was run on the 8,786 unresolved corpus rows. Section~\ref{sec:challenge} tests it on held-out posts.

\section{Evaluation Design and Evidence}

\subsection{Required metrics}

For a three-class sentiment task, overall accuracy alone can hide poor performance on minority classes. The evaluation module therefore calculates a confusion matrix and per-class precision, recall and F1, together with macro F1. For issue extraction, the appropriate measures are per-class precision, recall and F1 under a multi-label protocol, plus micro and macro averages.

Model comparison uses the same independently annotated records. Runtime is measured on the same machine, and the API run records input tokens, output tokens, retries, the exact returned model identifier and estimated cost from the provider list prices supplied at run time.

\subsection{Independent annotation}

A blinded evaluation package was generated without VADER, transformer or LLM columns. Threads were split so that representative and challenge posts do not share threads, with at most five comments per thread. Annotator~A labelled:

\begin{itemize}
  \item a 300-row representative sheet (214 Reddit, 86 HWZ), of which 295 posts were telco-relevant;
  \item a 100-row challenge sheet reserved for harder cases (negation, sarcasm or Singlish cues, mixed brands, weak VADER scores, long posts).
\end{itemize}

A second annotator labelled a 100-row overlap (60 representative, 40 challenge; \texttt{double\_annotation\_annotator\_b\_filled2.csv}). On the representative overlap, sentiment agreement is 62\% (Cohen's $\kappa=0.39$ \cite{cohen1960kappa}, $n=58$), which is fair agreement; on the challenge overlap it is 73\% ($\kappa=0.59$, $n=40$). Most disagreements share one pattern: posts Annotator~A labelled \texttt{not\_happy} that Annotator~B labelled \texttt{neutral} (9 of 22 disagreements on the representative overlap, 8 of 11 on the challenge overlap). They are typically questions about compensation, troubleshooting advice or second-hand reports. The annotation guide originally listed ``service failure'' under \texttt{not\_happy} and ``factual statement'' under \texttt{neutral} without saying which takes precedence; it now states that a first-hand failure report is \texttt{not\_happy}, and Annotator~B's relabelling applies that rule, while Annotator~A's labels predate it. Two earlier versions of Annotator~B's sheet are superseded: the first because its sentiment labels contradicted B's own notes and matched Annotator~A, the second because only sentiment had been revised and the other fields still matched Annotator~A's. The version used here was relabelled with B's own notes throughout; on the challenge overlap 9 of 40 rows match Annotator~A on every field and no notes are identical. Headline metrics use Annotator~A's representative labels; challenge metrics are reported separately and are not pooled.

An earlier 150-row file reported 100\% VADER accuracy. Those labels match VADER exactly and can be created by prefilling from the model, so that file is treated as leakage and is not used as evidence.

\subsection{Headline comparison (representative, $n=295$)}

Table~\ref{tab:headlinemetrics} summarises GPT-4o-mini and GPT-4o against the local baselines on the same 295 human-labelled posts. Sentiment support is 53 happy, 147 not-happy and 95 neutral. The hybrid issue row keeps keyword labels except where rules output only \texttt{general\_discussion}, then substitutes the GPT-4o issue list from the same run.

\begin{table}[H]
\centering
\caption{Representative-set comparison against Annotator~A labels}
\label{tab:headlinemetrics}
\begin{tabular}{lrrrr}
\toprule
\textbf{Task / method} & \textbf{Accuracy} & \textbf{Macro F1} & \textbf{Micro F1} & \textbf{Exact match} \\
\midrule
Sentiment: VADER & 0.366 & 0.360 & --- & --- \\
Sentiment: GPT-4o-mini & 0.671 & \textbf{0.631} & --- & --- \\
Sentiment: GPT-4o & 0.675 & 0.623 & --- & --- \\
Issues: keyword rules & --- & 0.639 & 0.657 & 0.380 \\
Issues: GPT-4o-mini & --- & 0.504 & 0.526 & 0.302 \\
Issues: GPT-4o (all rows) & --- & 0.600 & 0.637 & 0.346 \\
Issues: GPT-4o hybrid (unresolved only) & --- & \textbf{0.680} & \textbf{0.711} & \textbf{0.403} \\
\bottomrule
\end{tabular}
\end{table}

GPT-4o-mini matches GPT-4o on sentiment at about one-seventeenth of the cost. No all-row LLM beats keyword rules on issues: GPT-4o closes most of the mini gap (0.504 to 0.600) but does not exceed rules (0.639). The unresolved-only hybrid is the only issue method that beats rules on this set, because it keeps Devices and other keyword-heavy labels and spends the model only on the 98 gold rows that rules could not classify.

Paired bootstrap intervals (2,000 resamples of the 295 posts) separate real gains from ties. GPT-4o-mini beats VADER on sentiment by 0.271 macro F1 (95\% CI 0.191 to 0.347). GPT-4o-mini and GPT-4o are tied (difference 0.008, CI $-0.043$ to 0.064). On issues, rules and all-row GPT-4o are also tied (GPT-4o minus rules $-0.039$, CI $-0.096$ to 0.019), while the hybrid's gain over rules is small but reliable (0.041, CI 0.011 to 0.072). A second GPT-4o-mini run on the same 295 posts reproduced 292 of 295 sentiment labels (run-to-run $\kappa=0.98$; macro F1 0.631 and 0.627), so the model's own run-to-run variation is far smaller than these differences.

\begin{table}[H]
\centering
\caption{VADER sentiment confusion matrix, representative set (rows = human label)}
\label{tab:cmvader}
\begin{tabular}{lrrr}
\toprule
 & \textbf{Pred.\ happy} & \textbf{Pred.\ not happy} & \textbf{Pred.\ neutral} \\
\midrule
Actual happy & 31 & 10 & 12 \\
Actual not happy & 61 & 50 & 36 \\
Actual neutral & 45 & 23 & 27 \\
\bottomrule
\end{tabular}
\end{table}

\begin{table}[H]
\centering
\caption{GPT-4o-mini sentiment confusion matrix, representative set (rows = human label)}
\label{tab:cmllm}
\begin{tabular}{lrrr}
\toprule
 & \textbf{Pred.\ happy} & \textbf{Pred.\ not happy} & \textbf{Pred.\ neutral} \\
\midrule
Actual happy & 27 & 14 & 12 \\
Actual not happy & 4 & 116 & 27 \\
Actual neutral & 16 & 24 & 55 \\
\bottomrule
\end{tabular}
\end{table}

VADER assigns happy to 61 of 147 human not-happy posts. GPT-4o-mini recovers most complaints (recall 0.79) but still confuses happy with the other two classes (F1 0.54 on 53 examples). GPT-4o is slightly more conservative on happy (recall 0.38, F1 0.49) and does not improve macro F1 over mini.

\begin{table}[H]
\centering
\caption{Per-class sentiment metrics on the representative set}
\label{tab:sentclass}
\begin{tabular}{llrrrr}
\toprule
\textbf{Method} & \textbf{Class} & \textbf{Precision} & \textbf{Recall} & \textbf{F1} & \textbf{Support} \\
\midrule
VADER & Happy & 0.226 & 0.585 & 0.326 & 53 \\
VADER & Not happy & 0.602 & 0.340 & 0.435 & 147 \\
VADER & Neutral & 0.360 & 0.284 & 0.318 & 95 \\
GPT-4o-mini & Happy & 0.574 & 0.509 & 0.540 & 53 \\
GPT-4o-mini & Not happy & 0.753 & 0.789 & 0.771 & 147 \\
GPT-4o-mini & Neutral & 0.585 & 0.579 & 0.582 & 95 \\
GPT-4o & Happy & 0.690 & 0.377 & 0.488 & 53 \\
GPT-4o & Not happy & 0.767 & 0.782 & 0.774 & 147 \\
GPT-4o & Neutral & 0.552 & 0.674 & 0.607 & 95 \\
\bottomrule
\end{tabular}
\end{table}

On issues, all-row LLMs are often precise but incomplete: they return one or two topics when the annotator marked several. Devices is the clearest mini miss (recall 0.058; 3 of 52). GPT-4o raises Devices recall only to 0.25 (13 of 52), still far below rules (0.81). The hybrid keeps those rule hits and uses GPT-4o on the residual: of 98 gold rows that rules labelled \texttt{general\_discussion}, GPT-4o reassigned 70. Network, porting and billing improve; Devices F1 stays with the rules (0.761). Rules still over-assign general discussion when used alone (66 false positives).

\begin{table}[H]
\centering
\caption{Selected per-class issue F1 scores on the representative set}
\label{tab:issueclass}
\begin{tabular}{lrrrrr}
\toprule
\textbf{Label} & \textbf{Support} & \textbf{Rules} & \textbf{Mini} & \textbf{GPT-4o} & \textbf{Hybrid} \\
\midrule
Network quality & 154 & 0.735 & 0.672 & 0.796 & 0.788 \\
Mobile plans & 110 & 0.716 & 0.503 & 0.611 & 0.758 \\
Devices & 52 & 0.750 & 0.109 & 0.361 & 0.761 \\
App \& account & 44 & 0.627 & 0.456 & 0.475 & 0.647 \\
Roaming & 22 & 0.762 & 0.533 & 0.718 & 0.800 \\
Porting & 20 & 0.581 & 0.750 & 0.703 & 0.667 \\
General discussion & 40 & 0.464 & 0.446 & 0.608 & 0.559 \\
\bottomrule
\end{tabular}
\end{table}

\subsection{GPT-4o unresolved hybrid on the corpus}

Because the hybrid beat rules on gold labels, \texttt{processing/issue\_extraction.py} was run with \texttt{--use-llm --llm-unresolved-only --llm-model gpt-4o} on all 24,690 labelled rows, of which 12,497 had no rule match. The API snapshot was the same \textbf{gpt-4o-2024-08-06} used in evaluation, in batches of ten, with a resumable checkpoint. No batch failed. Keyword-only labels were kept in a separate file; the dashboard now overlays the hybrid labels.

\begin{table}[H]
\centering
\caption{Full-corpus issue counts: keyword rules versus GPT-4o on unresolved rows}
\label{tab:corpushybrid}
\begin{tabular}{lrrr}
\toprule
\textbf{Label} & \textbf{Rules} & \textbf{Hybrid} & \textbf{Change} \\
\midrule
General discussion & 3,229 & 1,342 & $-$1,887 \\
Network quality & 2,769 & 3,764 & $+$995 \\
Mobile plans & 2,079 & 2,376 & $+$297 \\
Pricing \& promotions & 901 & 1,144 & $+$243 \\
Customer service & 579 & 748 & $+$169 \\
Roaming & 853 & 959 & $+$106 \\
Devices & 1,611 & 1,691 & $+$80 \\
Porting & 527 & 600 & $+$73 \\
App \& account & 641 & 713 & $+$72 \\
Billing \& payment & 659 & 730 & $+$71 \\
\bottomrule
\end{tabular}
\end{table}

Counts are posts carrying each label in the 9,162-post clean corpus. General discussion fell from 35\% of posts to 15\%, and GPT-4o mainly pulled network talk out of the leftover bucket. Devices barely moved, which matches the gold set: rules already catch most device posts. These corpus counts have no human labels, so they are a production overlay rather than a second accuracy table.

\subsection{Held-out challenge set ($n=100$)}
\label{sec:challenge}

The challenge set was sampled from threads not used in the representative set and weighted toward hard cases (negation, sarcasm, Singlish, mixed brands, long posts). It is reported separately and never pooled. Because the hybrid was chosen after seeing representative-set results, the challenge set is its held-out test. GPT-4o was run on it for this purpose (\$0.050), and the hybrid was scored from those predictions without new API calls.

\begin{table}[H]
\centering
\caption{Challenge-set comparison against Annotator~A labels}
\label{tab:challenge}
\begin{tabular}{lrr}
\toprule
\textbf{Method} & \textbf{Sentiment macro F1} & \textbf{Issue macro F1} \\
\midrule
VADER / keyword rules & 0.491 & 0.499 \\
GPT-4o-mini (all posts) & \textbf{0.773} & 0.470 \\
GPT-4o (all posts) & 0.712 & \textbf{0.596} \\
Hybrid: rules, then GPT-4o on unresolved posts & --- & 0.509 \\
\bottomrule
\end{tabular}
\end{table}

The sentiment conclusion holds on hard posts. GPT-4o-mini beats VADER (difference 0.282, 95\% CI 0.138 to 0.413) and is not beaten by GPT-4o (mini minus GPT-4o 0.062, CI $-0.002$ to 0.127). The issue conclusion does not fully hold. The hybrid's gain over rules shrinks to 0.010 (CI $-0.039$ to 0.065), and GPT-4o labelling every post is best, beating both rules (0.097, CI 0.014 to 0.181) and the hybrid (0.087, CI 0.017 to 0.151). On these posts GPT-4o does better on customer service, roaming and billing, while rules still recall more device posts (0.75 versus 0.33). The hybrid remains the default for the corpus, because the representative set rather than the challenge set reflects its mix of posts, but its 0.680 should be read as optimistic. A union of rule and GPT-4o labels scored 0.703 on the representative set but 0.533 on the challenge set; it was found after seeing both results and is not adopted.

\subsection{Recent-posts gold set ($n=97$)}
\label{sec:recentgold}

The representative and challenge sets were drawn from posts up to early April~2026, so a third blind 100-post set was drawn from posts published between 10~April and 26~September~2026, with the same rules as the representative set (\texttt{processing/create\_recent\_gold\_set.py}: posts naming a telco, one copy per text, at most five per thread, 56 Reddit and 44 HWZ from 73 threads) and labelled by Annotator~A; 97 were telco-relevant. \texttt{models/score\_recent\_gold.py} scores the production labels already stored in the corpus, so no new API calls were made. Two misspelt issue labels in the annotation file were normalised at scoring time, and four relevant posts without an issue label were left out of the issue scores.

\begin{table}[H]
\centering
\caption{Recent-posts gold set: production labels against Annotator~A}
\label{tab:recentgold}
\begin{tabular}{lrr}
\toprule
\textbf{Method} & \textbf{Sentiment macro F1} & \textbf{Issue macro F1} \\
\midrule
VADER / keyword rules & 0.428 & 0.413 \\
GPT-4o-mini (all posts) & \textbf{0.637} & 0.440 \\
Hybrid: rules, then GPT-4o on unresolved posts & --- & \textbf{0.477} \\
\bottomrule
\end{tabular}
\end{table}

The sentiment result carries over: GPT-4o-mini beats VADER by 0.209 macro F1 (95\% CI 0.086 to 0.328), on both Reddit (0.651) and HWZ (0.614), and its score matches the representative set (0.631). The error is not symmetric, however. Of 39 posts the annotator labelled not happy, GPT-4o-mini labelled 15 neutral (not-happy recall 0.59), and it labels 41\% of the sample neutral against the annotator's 30\%. Corpus not-happy shares are therefore more likely understated than overstated. Issue tagging is harder on these recent posts for every method; the hybrid still beats rules (0.064, CI 0.023 to 0.107), but Customer Service (F1 0.24) and Devices (0.16, on two posts) are weak. With 97 posts and one annotator, per-class values are indicative only.

\subsection{Qualitative disagreements}

Examples are paraphrased from \texttt{qualitative\_disagreements.csv}. They illustrate why surface lexicons fail and why all-row LLMs drop device labels.

\begin{itemize}
  \item \textbf{REP-0003} (``SingTel eSIM still down''): human and LLM \texttt{not\_happy}; VADER \texttt{neutral}. Short status updates lack sentiment words.
  \item \textbf{REP-0014} (switched after 5G drops, cheaper 500\,GB plan): human and LLM \texttt{not\_happy}; VADER \texttt{happy} because of ``better'' and ``just \$12''.
  \item \textbf{CHL-0005} (``Smoooooooth as baby butttttt''): human and LLM \texttt{happy}; VADER \texttt{neutral}. Elongated slang is outside the lexicon.
  \item \textbf{CHL-0029} (nationwide outage, ready to port out): human and LLM \texttt{not\_happy}; VADER \texttt{happy}. Switching intent is a complaint, not praise.
  \item \textbf{REP-0015} (``Starhub, no ONR, no CGNAT''): human and rules label \texttt{devices}; GPT-4o-mini returns \texttt{general\_discussion} and flips sentiment to happy. Equipment jargon is a systematic all-row LLM miss; the unresolved hybrid keeps the rule \texttt{devices} tag on this class of post.
\end{itemize}

\subsection{Runtime and cost}

Persisted logs from \texttt{models/llm\_evaluation.py} and the unresolved corpus run:

\begin{table}[H]
\centering
\caption{Measured LLM cost and runtime}
\label{tab:cost}
\begin{tabular}{llrrrrr}
\toprule
\textbf{Job} & \textbf{Model} & \textbf{$n$} & \textbf{Seconds} & \textbf{Tokens} & \textbf{Est.\ USD} & \textbf{Fail} \\
\midrule
Representative eval & GPT-4o-mini & 295 & 74 & 38,464 & 0.0087 & 0 \\
Challenge eval & GPT-4o-mini & 100 & 28 & 12,948 & 0.0030 & 0 \\
Representative eval & GPT-4o & 295 & 73 & 38,545 & 0.145 & 0 \\
Challenge eval & GPT-4o & 100 & 20 & 12,920 & 0.050 & 0 \\
Representative rerun & GPT-4o-mini & 295 & 66 & 38,476 & 0.0087 & 0 \\
Corpus unresolved issues & GPT-4o & 12,497 & 2,445 & 1,269,218 & 4.80 & 0 \\
Corpus sentiment (posts to 9 Apr) & GPT-4o-mini & 1,711 & 382 & 204,623 & 0.048 & 0 \\
Corpus sentiment (all posts) & GPT-4o-mini & 9,162 & --- & 1,174,839 & 0.28 & 1 \\
\bottomrule
\end{tabular}
\end{table}

Mini prices were \$0.15 per million input tokens and \$0.60 per million output tokens. GPT-4o prices were \$2.50 / \$10.00. VADER labelled the representative set in 0.22 seconds and keyword rules in 0.27 seconds, with no API charge. Twitter-RoBERTa was not run because \texttt{transformers} was not installed. The hybrid gold metrics in Tables~\ref{tab:headlinemetrics} and~\ref{tab:challenge} reuse the GPT-4o predictions and added no API cost. Per post, GPT-4o costs about \$0.0005 and GPT-4o-mini about \$0.00003. A first pass on the April--September posts that sent Reddit posts without thread titles (about \$1.10 of GPT-4o and \$0.05 of GPT-4o-mini) was discarded and is not included above.

\section{Findings}

RQ3 is answered in Section~4: GPT-4o-mini adds sentiment context over VADER on typical and hard posts; for issues, rules plus GPT-4o on unresolved posts is best on typical posts, and GPT-4o labelling every post is best on hard ones. This section applies GPT-4o-mini sentiment and hybrid issue labels to the corpus: 9,162 unique posts in 349 telco threads, collected up to 26~September~2026. Shares are reported per post and with every thread weighted equally, so that one large thread cannot drive a finding. Section~\ref{sec:casestudy} then shows how one outage distorts these aggregates.

\subsection{Sentiment: most posts are negative or neutral}

\begin{table}[H]
\centering
\caption{Sentiment shares by labeller}
\label{tab:corpussentiment}
\begin{tabular}{lrrrr}
\toprule
\textbf{Labeller and data} & \textbf{Happy} & \textbf{Not happy} & \textbf{Neutral} & \textbf{$n$} \\
\midrule
GPT-4o-mini, per post & 19\% & 44\% & 37\% & 9,162 \\
GPT-4o-mini, thread-weighted & 18\% & 42\% & 40\% & 349 threads \\
GPT-4o-mini, outside outage threads & 20\% & 40\% & 40\% & 8,060 \\
VADER, same posts & 46\% & 28\% & 27\% & 9,162 \\
Human labels, representative gold set & 18\% & 50\% & 32\% & 295 \\
\bottomrule
\end{tabular}
\end{table}

About two in five posts are not happy under every weighting, and fewer than one in five is happy (RQ1). VADER again nearly reverses the picture: of 4,180 posts it calls happy, GPT-4o-mini reads 1,480 as not happy and 1,538 as neutral. The recent-posts gold set (Section~\ref{sec:recentgold}) shows the direction of GPT-4o-mini's error on April--September posts: it labels 41\% neutral where the human labels say 30\%, mainly by reading complaints as neutral, so the corpus not-happy share is more likely understated than overstated. These shares describe unsolicited forum discussion and are not a satisfaction rate.

\subsection{Issue themes}

\begin{table}[H]
\centering
\caption{Primary issue themes}
\label{tab:overallfindings}
\begin{tabular}{lrrrrr}
\toprule
\textbf{Issue} & \textbf{Posts} & \textbf{Threads} & \textbf{Share (post)} & \textbf{Share (thread)} & \textbf{Not happy} \\
\midrule
Network Quality & 3,733 & 267 & 40.7\% & 36.9\% & 53\% \\
General Discussion & 1,342 & 178 & 14.6\% & 13.0\% & 29\% \\
Mobile Plans & 1,263 & 202 & 13.8\% & 12.6\% & 37\% \\
Pricing \& Promotions & 745 & 163 & 8.1\% & 6.6\% & 42\% \\
Devices & 545 & 145 & 5.9\% & 6.5\% & 35\% \\
Customer Service & 338 & 115 & 3.7\% & 3.9\% & 59\% \\
App \& Account & 335 & 115 & 3.7\% & 9.2\% & 40\% \\
Roaming & 335 & 81 & 3.7\% & 4.9\% & 31\% \\
Billing \& Payment & 264 & 90 & 2.9\% & 3.5\% & 53\% \\
Porting & 262 & 58 & 2.9\% & 2.8\% & 35\% \\
\bottomrule
\end{tabular}
\end{table}

Network quality leads under both weightings and carries 46\% of the thread-weighted not-happy discussion. Mobile-plan comparisons are second and pricing third. App and account access is small per post (3.7\%) but appears in 115 threads, so with threads weighted equally it rises to 9.2\%: it is a steady low-volume complaint across many conversations rather than a few large ones. Customer service and billing are the most negative themes by rate (59\% and 53\% not happy) but small in volume.

\subsection{Platform differences}

Per post, Reddit is 48\% not happy and HWZ 37\%. Outside the outage threads the gap narrows to 42\% against 37\%, and human labels on the representative gold set show 51\% against 47\%. The remaining five points are confounded with topic: HWZ discussion is dominated by long plan and MVNO megathreads, which carry more neutral comparison, while Reddit carries more incident reports and scam warnings. For RQ2, the platforms differ mainly in what they discuss; the data do not support a claim that one platform's users are more negative.

\subsection{Cross-brand observations}

\begin{table}[H]
\centering
\caption{Not-happy share for brands with at least 92 posts and 10 threads}
\label{tab:brands}
\begin{tabular}{lrrrrr}
\toprule
\textbf{Brand} & \textbf{Posts} & \textbf{Threads} & \textbf{Not happy} & \textbf{Thread-weighted} & \textbf{Outside outage threads} \\
\midrule
Singtel & 3,970 & 238 & 48\% & 45\% & 42\% \\
StarHub & 1,638 & 180 & 44\% & 40\% & 40\% \\
SIMBA & 1,461 & 152 & 44\% & 45\% & 43\% \\
M1 & 1,415 & 156 & 40\% & 43\% & 40\% \\
Eight & 734 & 109 & 35\% & 31\% & 35\% \\
Zym & 580 & 73 & 35\% & 42\% & 35\% \\
Giga & 360 & 69 & 46\% & 39\% & 43\% \\
\bottomrule
\end{tabular}
\end{table}

Outside outage threads, Singtel, StarHub, SIMBA, M1 and Giga sit within three points of each other (40--43\%). Eight is consistently the least negative (31--35\% under every weighting); Zym is low per post but not thread-weighted, so it is not a robust signal. No brand league table is defensible: the differences are small, thread-weighting reorders several brands, and each post carries one sentiment label copied to every brand it names. redONE (53 posts in 13 threads) is below the ranking threshold. Mention counts are not unique customers.

\subsection{Case study: the March 2026 Singtel outage}
\label{sec:casestudy}

Forty-five threads are about outages, led by ``16 March 26: Singtel mobile service outage''. They hold 1,102 posts (12\%), and they show how one incident distorts every aggregate. Inside them 71\% of posts are not happy, against 40\% in all other threads. Singtel is 73\% not happy inside them and 42\% outside, next to StarHub (68\% and 40\%). And 1,099 of the 1,102 outage posts are on Reddit, against 57\% of other posts, so the outage threads create most of Reddit's per-post lead over HWZ (48\% against 37\%, but 42\% against 37\% outside them). Many billing and customer-service posts in these threads are outage aftercare: whether compensation is a refund or a goodwill credit, and support lines that could not be reached.

The weekly huddle view catches the event. With a threshold set relative to a normal week (Section~\ref{sec:dashboard}), only two weeks qualify: the outage week (352 posts) and the week of 14~September~2026 (149 posts). The second is a different kind of event, spread across brands: StarHub broadband going down, OTP failures on the SIMBA app, and complaints in the SIMBA, Singtel, M1 Maxx and Eight threads. The screen therefore flags both a single-operator outage and a multi-brand bad week, which is the early-warning use the project set out to support.

\subsection{Dashboard design}
\label{sec:dashboard}

\texttt{dashboard/generate\_dashboard.py} builds a single static page, \texttt{dashboard/index.html}, from the clean corpus. Chart titles are computed from the data and state the takeaway (for example, ``Network Quality leads, then Mobile Plans''). Where a result would invite a wrong reading, the title says so: a per-post Reddit--HWZ gap is shown as ``check topic mix first'', not as a platform effect. The page opens with a weekly 30-minute huddle playbook and a weekly queue chart: (a) not-happy posts in Network Quality, App \& Account or Porting; (b) open ten source posts; (c) match them to tickets or outages; (d) assign one owner; (e) check next week whether the share fell. A fixed threshold does not scale with volume: a 20-post threshold is crossed in 27 weeks. The page therefore acts when a week reaches twice the median of the last 26 weeks (115 posts).

\begin{itemize}
  \item \textbf{Tone:} GPT-4o-mini shares beside VADER and the human gold labels, so readers see why VADER is not used.
  \item \textbf{Source:} Reddit and HWZ on the same labeller, with the gold-label comparison and the topic-mix caveat.
  \item \textbf{Issue volume:} per-post and thread-weighted shares side by side; a finding that holds in both is not one thread talking.
  \item \textbf{Not-happy by issue:} which themes carry the negative discussion.
  \item \textbf{Brands:} ranked only with at least 1\% of posts and 10 threads, with thread counts in the table.
  \item \textbf{Recommendations:} ordered by each issue's share of not-happy discussion.
\end{itemize}

Outputs remain screening signals, not verified customer facts.

\section{Business Recommendations}

Each recommendation names an owner, an action and a check. They are ordered by each issue's share of not-happy discussion in the corpus. Forum counts are a triage queue, not proof of market failure.

\subsection{Network: communicate during incidents, then prove local fixes}

Network quality carries 46\% of not-happy discussion. Of 1,991 not-happy network posts, 621 sit in outage threads, and the huddle flagged both the March outage and the September StarHub broadband disruption. The first action is incident communication: a status page and proactive updates with an estimated restoration time, posted where customers are already asking. Outside outages, 289 not-happy network posts name a place such as MRT tunnels, homes or offices; network teams should join those posts to tickets and cell data. Success is a shorter time to the first public update during an incident, and fewer tickets at a named location after remediation.

\subsection{Product and commercial: make plan terms comparable}

Mobile plans (13\% of posts, in 202 threads) and pricing (8\%, 42\% not happy) are the second and third sources of not-happy discussion. Posts centre on terms customers discover late: fair-use limits on ``unlimited'' plans, add-on charges such as Caller ID (which posts note the three network operators charge for while SIMBA includes it), introductory prices that later change, and which host network an MVNO uses. Product and commercial teams should show fair-use rules, add-on charges, recurring price and host network at checkout. Success is fewer ``which plan?'' and hidden-charge threads.

\subsection{Devices and digital: fix equipment and app friction}

Device posts (6\%) cover eSIM failures and ONR/bridge-mode router policy; device support should publish router and eSIM requirements per plan. App and account problems are small per post but spread across 115 threads, including OTP failures on the SIMBA app in the September spike week; digital teams should watch forum posts after each app release. Success is fewer equipment threads and fewer ``cannot log in'' posts after a release.

\subsection{Customer service and billing: plan the incident aftercare}

Customer service (59\% not happy) and billing (53\%) are small but the most negative themes, and in the outage case study they were largely aftercare. Service and billing teams should publish the compensation rule with the incident notice and keep a self-service path open during outages. Success is fewer compensation and ``cannot reach support'' posts after the next incident.

\subsection{Lower-volume themes: roaming and porting}

Roaming (4\%) and porting (3\%) are small and mostly questions (31\% and 35\% not happy), so a ``travelling to X'' checklist and clearer port timelines are low-cost fixes rather than priorities.

\subsection{Insight huddle: operate the dashboard as an alert-and-review tool}

Do not publish a brand league table. A weekly 30-minute process is enough:

\begin{enumerate}
  \item flag not-happy posts in Network, App \& Account or Porting only when that week's volume clears a threshold set relative to a normal week (twice the recent median);
  \item open about ten source records and discard sarcasm, mixed-brand and off-target labels;
  \item triangulate with tickets, outages and billing;
  \item assign one owner and a response date;
  \item next week, check whether that issue's share fell.
\end{enumerate}

Use GPT-4o-mini for sentiment on the whole feed: labelling all 9,162 posts cost about \$0.28, so a daily run costs cents. For issues, keep keyword rules and apply GPT-4o to \texttt{general\_discussion} rows. On hard posts, GPT-4o labelling every post is more accurate and, at about \$0.0005 per post, affordable on a deduplicated feed. Keep a person in the loop before any public reply. That mix is the business reading of the evaluation, not a model contest.

\section{Limitations and Ethical Considerations}

\subsection{Representativeness and platform bias}

Forum contributors are self-selecting and are not Singapore mobile users at large. Findings use a corpus of 9,162 unique posts, after removing repeated copies (24,690 raw rows held 11,643 unique posts) and 96 off-topic threads. The corpus is 62\% Reddit, and 95\% of posts are from 2026; the few older posts are opening pages of long HWZ threads. Twelve per cent of posts sit in outage threads. Two collection artefacts are worth stating. Early HWZ scrapes read only the first page of each thread because of a page-counting bug; later scrapes add the newest pages. Reddit posts after 9~April~2026 come from the Arctic Shift research archive rather than Reddit directly, because Reddit blocks unauthenticated access and gates new API applications; the archive can miss content deleted before capture, and it includes replies that the earlier Reddit scraper skipped. There are no app-store reviews, tickets or NPS in the analysed set, so results cannot be generalised to market share, population sentiment, churn or Net Promoter Score.

\subsection{Unequal and dependent observations}

One record can mention several brands, so mention counts are not unique customers. Comments in a thread are not independent: they respond to a shared outage or launch. The report therefore shows thread-weighted shares beside per-post shares, and brand shares outside outage threads. Mixed-brand posts were often labelled \texttt{neutral}, which can hide a Singtel complaint sitting next to StarHub praise.

\subsection{Label quality and LLM misclassification}

GPT-4o-mini matches Annotator~A on about two-thirds of posts (accuracy 0.67, macro F1 0.631), and happy is its weakest class (F1 0.54). The benchmark is itself one annotator: on the representative overlap two humans agree only at $\kappa=0.39$, mostly because the guide did not say whether a bare failure report is \texttt{not\_happy} or \texttt{neutral}. Annotator~A treats such reports as complaints, and GPT-4o-mini agrees with A on the examples in Section~4, so part of its margin over VADER depends on that reading. On the recent-posts gold set, sentiment accuracy holds (macro F1 0.637) but GPT-4o-mini recalls only 59\% of human not-happy posts, reading many as neutral, and issue tagging is markedly weaker than on the original posts (hybrid macro F1 0.477 against 0.680); that sample is 97 posts, labelled by Annotator~A alone, with two spelling fixes applied at scoring time. One of 9,162 posts could not be labelled by GPT-4o-mini and falls back to VADER; 40 posts that failed in batches of ten were relabelled one at a time. A first pass on the April--September posts sent Reddit posts without their thread titles; those labels were discarded and the posts relabelled with titles. All-post LLMs under-label Devices relative to rules (GPT-4o recall 0.25 versus 0.81 on the representative set, 0.33 versus 0.75 on the challenge set) and often return one topic when the post has several. Keyword rules over-trigger on generic terms such as ``plan''. The hybrid's 0.680 was measured on the set used to choose it; on held-out hard posts its gain over rules was not significant and GPT-4o on every post did better. The corpus issue reassignments have no second human check. These errors are prompt- and snapshot-specific. Forum text sent to a vendor API is also a privacy limit: production jobs must not submit NRIC, phone numbers or other identifiers. The defensible stance is triage plus human review, not automated league tables.

\subsection{Privacy and responsible use}

Although forum posts are publicly accessible, public availability does not remove ethical obligations. Reports should avoid usernames and unnecessary verbatim quotations, minimise stored identifiers, restrict raw-data access and comply with platform terms. The system should analyse themes, not profile individuals. Credentials and API keys must be stored in environment variables or a secret manager and must never appear in code, reports or version control.

\section{Future Work}

\begin{enumerate}
  \item Improve issue tagging on recent posts, which scores well below the representative set (hybrid macro F1 0.477), starting with Customer Service (F1 0.24), and have a second annotator label part of the recent-posts gold set.
  \item Adjudicate the remaining disagreements between the two annotators in a separate file, and relabel Annotator~A's overlap rows under the new guide rule for failure reports.
  \item Compare the hybrid with GPT-4o labelling every post on a fresh held-out sample, since the two gold sets disagree on which is better.
  \item Spot-check the corpus rows that GPT-4o moved out of \texttt{general\_discussion}, especially the extra network-quality tags, before treating huddle counts as incidents.
  \item Add target-aware sentiment, so that a post praising one operator and criticising another yields a label per brand.
  \item Validate dashboard alerts against internal service tickets and network events before operational deployment.
\end{enumerate}

\section{Conclusion}

This project demonstrates a complete technical path from online forum collection to operator action. On 9,162 unique posts across six months, network quality dominates both the volume and the negativity of discussion, followed by plan terms and pricing. About two in five posts are not happy under every weighting; Reddit and HWZ differ mainly in topic rather than tone, and no brand stands out beyond small, weighting-sensitive differences, with Eight consistently the least negative. The original, outage-heavy corpus shows why this matters: one incident inflated negativity, made Singtel look uniquely negative and made Reddit look more negative than HWZ. The weekly huddle view flagged that outage and a later multi-brand bad week, and nothing else.

The analysis also shows why model governance matters. A perfect score from non-independent labels was discarded in favour of an imperfect but credible benchmark. On 295 blinded human labels, gpt-4o-mini-2024-07-18 raises sentiment macro F1 from 0.360 (VADER) to 0.631; the gain holds on hard posts (0.773 versus 0.491), gpt-4o-2024-08-06 adds nothing for sentiment at about 17 times the cost, and repeated runs are stable. For issues the evidence is split: rules plus GPT-4o on unresolved posts is best on typical posts (0.680), and GPT-4o labelling every post is best on hard ones (0.596). The defensible deployment is LLM sentiment, rules plus a targeted LLM pass for issues, and human review of alerts, as an early-warning layer alongside, rather than in place of, customer-service and network data.

\section*{Reproducibility Note}
\addcontentsline{toc}{section}{Reproducibility Note}

From the repository root, the main workflow is:

\begin{verbatim}
python processing/preprocess.py
python processing/telco_labeler.py
python processing/issue_extraction.py
# Keyword-only labels. Hybrid GPT-4o unresolved run is already saved as
# data/processed/issues_posts_labeled_2026-08-21_203752_gpt4o_hybrid.csv
python processing/build_clean_corpus.py     # posts up to 2026-04-09

# Posts from 2026-04-10 (newest HWZ pages; Reddit via the Arctic Shift archive)
python scrapers/hwz_forum_scraper.py --forum-pages 5 --max-threads 60 \
  --max-thread-pages 15 --latest-pages --since 2026-04-10 --out data/raw/apr_sep_2026/hwz_latest
python scrapers/reddit_archive_scraper.py --subreddits singapore,askSingapore \
  --since 2026-04-10 --out data/raw/apr_sep_2026/reddit_archive
python processing/prepare_new_posts.py          # clean, label, keyword issues
# then the GPT-4o hybrid pass printed by prepare_new_posts.py, and:
python processing/prepare_new_posts.py --merge  # -> data/processed/corpus_posts.csv
python models/corpus_llm_sentiment.py --in data/processed/corpus_posts.csv \
  --out data/processed/corpus_posts_mini_sentiment.csv \
  --summary-out data/processed/corpus_posts_mini_sentiment_summary.json
python dashboard/generate_dashboard.py      # -> dashboard/index.html
python processing/build_gold_labels.py      # three gold sets -> gold_labels.csv
python models/score_recent_gold.py          # recent-posts gold set; no API calls
python models/robustness_checks.py          # CIs, stability, kappa; no API calls

# Independent LLM vs baseline evaluation (after human labels)
python models/llm_evaluation.py \
  --in data/evaluation/representative_annotator_a.csv \
  --out-dir data/evaluation/llm_comparison_representative \
  --model gpt-4o-mini \
  --input-cost-per-million 0.15 \
  --output-cost-per-million 0.60
python models/llm_evaluation.py \
  --in data/evaluation/challenge_annotator_a.csv \
  --out-dir data/evaluation/llm_comparison_challenge \
  --model gpt-4o-mini \
  --input-cost-per-million 0.15 \
  --output-cost-per-million 0.60
python models/llm_evaluation.py \
  --in data/evaluation/representative_annotator_a.csv \
  --out-dir data/evaluation/llm_comparison_representative_gpt4o \
  --model gpt-4o \
  --input-cost-per-million 2.50 \
  --output-cost-per-million 10.00
python models/llm_evaluation.py \
  --in data/evaluation/challenge_annotator_a.csv \
  --out-dir data/evaluation/llm_comparison_challenge_gpt4o \
  --model gpt-4o \
  --input-cost-per-million 2.50 \
  --output-cost-per-million 10.00
\end{verbatim}

Completed GPT-4o issue extraction on unresolved rows (already run; outputs persist):

\begin{verbatim}
python processing/issue_extraction.py \
  --in data/processed/posts_labeled_2026-08-21_203752.csv \
  --out data/processed/issues_posts_labeled_2026-08-21_203752_gpt4o_hybrid.csv \
  --summary-out data/processed/issues_summary_gpt4o_hybrid.json \
  --use-llm --llm-unresolved-only --llm-model gpt-4o \
  --input-cost-per-million 2.50 \
  --output-cost-per-million 10.00
\end{verbatim}

API credentials should be supplied through environment variables and excluded from all submitted material.

\begin{thebibliography}{9}

\bibitem{hutto2014vader}
C.~J. Hutto and E.~Gilbert,
``VADER: A Parsimonious Rule-based Model for Sentiment Analysis of Social Media Text,''
\textit{Proceedings of the Eighth International AAAI Conference on Weblogs and Social Media},
2014.

\bibitem{loureiro2022tweeteval}
D.~Loureiro, F.~Barbieri, L.~Neves, L.~E. Anke and J.~Camacho-Collados,
``TimeLMs: Diachronic Language Models from Twitter,''
\textit{Proceedings of the 60th Annual Meeting of the Association for Computational Linguistics: System Demonstrations},
pp.~251--260, 2022.

\bibitem{vaswani2017attention}
A.~Vaswani et al.,
``Attention Is All You Need,''
\textit{Advances in Neural Information Processing Systems}, vol.~30, 2017.

\bibitem{cohen1960kappa}
J.~Cohen,
``A Coefficient of Agreement for Nominal Scales,''
\textit{Educational and Psychological Measurement}, vol.~20, no.~1, pp.~37--46, 1960.

\end{thebibliography}

\appendix
\section{Repository-to-Report Traceability}

\begin{longtable}{@{}p{0.33\textwidth}p{0.61\textwidth}@{}}
\toprule
\textbf{Repository artifact} & \textbf{Evidence used in this report} \\
\midrule
\endfirsthead
\toprule
\textbf{Repository artifact} & \textbf{Evidence used in this report} \\
\midrule
\endhead
\texttt{README.md} & Pipeline scope, command sequence, script responsibilities and default models \\
\texttt{Project\_Progress.md} & Data inventory, stated milestones, dashboard findings and planned evaluation work \\
\texttt{processing/preprocess.py} & Normalisation and filtering design \\
\texttt{processing/telco\_labeler.py} & Brand matching, VADER thresholds and optional GPT-4o-mini sentiment \\
\texttt{processing/issue\_extraction.py} & Taxonomy, keyword rules and batched GPT-4o unresolved hybrid \\
\texttt{models/llm\_evaluation.py} & Blind comparison of VADER, keyword rules, GPT-4o-mini and GPT-4o \\
\texttt{llm\_comparison\_representative/metrics.json} & Mini $n=295$ confusion matrices, per-class P/R/F1, tokens, runtime and cost \\
\texttt{llm\_comparison\_representative\_gpt4o/metrics.json} & GPT-4o $n=295$ metrics, tokens, runtime and cost \\
\texttt{llm\_comparison\_representative\_gpt4o\_hybrid/metrics.json} & Unresolved-only hybrid scored on the same gold labels \\
\texttt{llm\_comparison\_challenge/metrics.json} & Separate $n=100$ hard-case mini metrics \\
\texttt{llm\_comparison\_challenge\_gpt4o/} and \texttt{\_hybrid/} & GPT-4o and hybrid on the held-out challenge set \\
\texttt{llm\_comparison\_representative\_rerun/} & Second GPT-4o-mini run for stability \\
\texttt{robustness\_summary.json} & Bootstrap intervals, stability, $\kappa$ per set, corpus outage and brand checks \\
\texttt{processing/build\_clean\_corpus.py}, \texttt{corpus\_summary.json} & Deduplication, off-topic thread filter, per-thread shares \\
\texttt{models/corpus\_llm\_sentiment.py}, \texttt{corpus\_posts\_mini\_sentiment\_summary.json} & Corpus GPT-4o-mini sentiment, tokens and cost \\
\texttt{qualitative\_disagreements.csv} & Rows where methods disagree with human labels \\
\texttt{sentiment\_benchmark.json} & Earlier leaked VADER 100\% file; not used as evidence \\
\texttt{issues\_summary.json} & Keyword-only full-corpus issue-label counts \\
\texttt{issues\_summary\_gpt4o\_hybrid.json} & Hybrid full-corpus issue-label counts after GPT-4o on unresolved rows \\
\texttt{dashboard/generate\_dashboard.py} & Dashboard (\texttt{index.html}): computed chart titles, weekly queue, recommendation order \\
\texttt{scrapers/reddit\_archive\_scraper.py}, \texttt{hwz\_forum\_scraper.py} & Collection from 10~April~2026 (Arctic Shift archive; HWZ newest pages) \\
\texttt{processing/prepare\_new\_posts.py} & Cleaning, labelling and issue rules for posts from 10~April, and merge into the full corpus \\
\texttt{processing/create\_recent\_gold\_set.py}, \texttt{recent\_sampling\_metadata.json} & Blind recent-posts gold sample \\
\texttt{processing/build\_gold\_labels.py}, \texttt{gold\_labels.csv} & The three gold sets in one file \\
\texttt{models/score\_recent\_gold.py}, \texttt{recent\_gold\_scores.json} & Production labels scored on the recent-posts gold set \\
\texttt{dashboard/Dash\_README.md} & Dashboard input row counts and generated summary \\
\bottomrule
\end{longtable}

\section{Evaluation Checklist Before Submission}

\begin{itemize}
  \item Replace the title-page placeholders with the official student, course and supervisor details.
  \item Do not describe the gold labels as dual-adjudicated; two annotators overlap on 100 posts but disagreements are not yet adjudicated.
  \item Confirm that no credentials are included in the submission archive, Terminal history files or Git history.
  \item Remove or anonymise usernames and any unnecessary personal identifiers.
\end{itemize}

\end{document}
